
\documentclass{article}
\usepackage{colortbl}
\usepackage{makecell}
\usepackage{multirow}
\usepackage{supertabular}

\begin{document}

\newcounter{utterance}

\centering \large Interaction Transcript for game `hot\_air\_balloon', experiment `air\_balloon\_survival\_lv\_complexity\_easy', episode 4 with Qwen3.6{-}35B{-}A3B{-}FP8{-}without{-}reasoning.
\vspace{24pt}

{ \footnotesize  \setcounter{utterance}{1}
\setlength{\tabcolsep}{0pt}
\begin{supertabular}{c@{$\;$}|p{.15\linewidth}@{}p{.15\linewidth}p{.15\linewidth}p{.15\linewidth}p{.15\linewidth}p{.15\linewidth}}
    \# & $\;$A & \multicolumn{4}{c}{Game Master} & $\;\:$B\\
    \hline

    \theutterance \stepcounter{utterance}  
    & & \multicolumn{4}{p{0.6\linewidth}}{
        \cellcolor[rgb]{0.9,0.9,0.9}{
            \makecell[{{p{\linewidth}}}]{
                \texttt{\tiny{[P1$\langle$GM]}}
                \texttt{Jūs piedalāties kopīgā sarunu spēlē.} \\
\\ 
\texttt{Kopā ar citu dalībnieku jums jāvienojas par vienu priekšmetu kopumu, kas tiks saglabāts. Katram no jums ir savs viedoklis par to, cik liela nozīme jums ir katram priekšmetam (nozīmīgums). Jūs nezināt, kā otrs dalībnieks vērtē priekšmetus. Turklāt jums ir norādīts, cik pūļu prasa katrs priekšmets.} \\
\texttt{Jūs varat vienoties par kopumu tikai tad, ja izvēlēto priekšmetu kopējais piepūles apjoms nepārsniedz kopīgo robežu:} \\
\\ 
\texttt{LIMIT = 4398} \\
\\ 
\texttt{Šeit ir norādītas atsevišķu priekšmetu piepūles vērtības:} \\
\\ 
\texttt{Item effort = \{"C08": 253, "C55": 370, "A43": 84, "A67": 840, "B42": 288, "A62": 92, "B19": 772, "C43": 459, "C73": 93, "C10": 668, "B36": 589, "A18": 659, "C61": 348, "C67": 964, "C42": 233, "C70": 400, "B40": 990, "C39": 315, "A21": 43, "C16": 336\}} \\
\\ 
\texttt{Šeit ir jūsu personīgais viedoklis par katra priekšmeta nozīmīgumu:} \\
\\ 
\texttt{Item importance values = \{"C08": 138, "C55": 583, "A43": 868, "A67": 822, "B42": 783, "A62": 65, "B19": 262, "C43": 121, "C73": 508, "C10": 780, "B36": 461, "A18": 484, "C61": 668, "C67": 389, "C42": 808, "C70": 215, "B40": 97, "C39": 500, "A21": 30, "C16": 915\}} \\
\\ 
\texttt{Mērķis:} \\
\\ 
\texttt{Jūsu mērķis ir vienoties par tādu kopīgu priekšmetu kopumu, no kura Jūs gūsiet pēc iespējas lielāku labumu (t.i., kas maksimāli palielina kopējo nozīmīgumu JUMS), vienlaikus nepārsniedzot LIMIT. Katrā ziņojumā jums nav obligāti jāiesniedz PRIEKŠLIKUMS {-} jūs varat vienkārši arī sarunāties. Visas taktikas ir atļautas!} \\
\\ 
\texttt{Saskarsmes protokols:} \\
\\ 
\texttt{Ziņojumā drīkst izmantot tikai šādus strukturētus formātus:} \\
\\ 
\texttt{PRIEKŠLIKUMS: \{'A', 'B', 'C', …\}} \\
\texttt{Ierosiniet saglabāt tieši šos priekšmetus.} \\
\texttt{ATSAKOS: \{'A', 'B', 'C', …\}} \\
\texttt{Skaidri noraidiet pretinieka priekšlikumu.} \\
\texttt{ARGUMENTS: \{'...'\}} \\
\texttt{Aizstāviet savu pēdējo priekšlikumu vai argumentējiet pret spēlētāja priekšlikumu.} \\
\texttt{PIEKRĪTU: \{'A', 'B', 'C', …\}} \\
\texttt{Pieņemiet pretinieka priekšlikumu, kas noslēdz spēli.} \\
\\ 
\\ 
\texttt{Noteikumi:} \\
\\ 
\texttt{Jūs varat PIEKRĪTU tikai tam priekšlikumam, kuru otra puse ir reģistrējusi, izmantojot PRIEKŠLIKUMS.} \\
\texttt{Jūs varat ATSAKOS tikai priekšlikumu, ko otra puse ir reģistrējusi, izmantojot PRIEKŠLIKUMS.} \\
\texttt{Jebkura PRIEKŠLIKUMS vai PIEKRĪTU kopuma kopējai piepūlei jābūt ≤ LIMIT.} \\
\texttt{NEATKLĀJIET savus slēptos nozīmīguma rādītājus.} \\
\texttt{Pēc strukturēta formāta taga ir jāseko kolu zīmei un atstarpei. Argumentam jābūt python kopai, kas satur 0 vai vairāk virknes.} \\
\texttt{Tātad tam jābūt formā "TAG: \{...\}"} \\
\texttt{Stingri ievērojiet saskarsmes protokolu un NERAKSTIET neko ārpus dotās struktūras.} \\
\texttt{Spēle beidzas, kad viena no pusēm dod PIEKRĪTU otra spēlētāja izteiktajam PRIEKŠLIKUMS.} \\
\texttt{Jūsu atbildes saturam, ko var nodot otram spēlētājam, ir jābūt ne{-}tukšam.} \\
\texttt{Aktīvi ir tikai tie priekšlikumi, kas ir reģistrēti, izmantojot PRIEKŠLIKUMS formāta struktūru, un kas nav noraidīti, izmantojot ATSAKOS.} \\
\\ 
            }
        }
    }
    & \\ \\

    \theutterance \stepcounter{utterance}  
    & \multicolumn{4}{p{0.6\linewidth}}{
        \cellcolor[rgb]{0.9,0.9,0.9}{
            \makecell[{{p{\linewidth}}}]{
                \texttt{\tiny{[P1$\rangle$GM]}}
                \texttt{PRIEKŠLIKUMS: \{"A21", "C73", "C43", "C61", "C39", "C16", "B42", "A43", "C55"\}} \\
            }
        }
    }
    & & \\ \\

    \theutterance \stepcounter{utterance}  
    & & \multicolumn{4}{p{0.6\linewidth}}{
        \cellcolor[rgb]{0.9,0.9,0.9}{
            \makecell[{{p{\linewidth}}}]{
                \texttt{\tiny{[GM$\rangle$P2]}}
                \texttt{Jūs piedalāties kopīgā sarunu spēlē.} \\
\\ 
\texttt{Kopā ar citu dalībnieku jums jāvienojas par vienu priekšmetu kopumu, kas tiks saglabāts. Katram no jums ir savs viedoklis par to, cik liela nozīme jums ir katram priekšmetam (nozīmīgums). Jūs nezināt, kā otrs dalībnieks vērtē priekšmetus. Turklāt jums ir norādīts, cik pūļu prasa katrs priekšmets.} \\
\texttt{Jūs varat vienoties par kopumu tikai tad, ja izvēlēto priekšmetu kopējais piepūles apjoms nepārsniedz kopīgo robežu:} \\
\\ 
\texttt{LIMIT = 4398} \\
\\ 
\texttt{Šeit ir norādītas atsevišķu priekšmetu piepūles vērtības:} \\
\\ 
\texttt{Item effort = \{"C08": 253, "C55": 370, "A43": 84, "A67": 840, "B42": 288, "A62": 92, "B19": 772, "C43": 459, "C73": 93, "C10": 668, "B36": 589, "A18": 659, "C61": 348, "C67": 964, "C42": 233, "C70": 400, "B40": 990, "C39": 315, "A21": 43, "C16": 336\}} \\
\\ 
\texttt{Šeit ir jūsu personīgais viedoklis par katra priekšmeta nozīmīgumu:} \\
\\ 
\texttt{Item importance values = \{"C08": 138, "C55": 583, "A43": 868, "A67": 822, "B42": 783, "A62": 65, "B19": 262, "C43": 121, "C73": 508, "C10": 780, "B36": 461, "A18": 484, "C61": 668, "C67": 389, "C42": 808, "C70": 215, "B40": 97, "C39": 500, "A21": 30, "C16": 915\}} \\
\\ 
\texttt{Mērķis:} \\
\\ 
\texttt{Jūsu mērķis ir vienoties par tādu kopīgu priekšmetu kopumu, no kura Jūs gūsiet pēc iespējas lielāku labumu (t. i., kas maksimāli palielina kopējo nozīmīgumu JUMS), vienlaikus nepārsniedzot LIMIT. Katrā ziņojumā jums nav obligāti jāiesniedz PRIEKŠLIKUMS {-} jūs varat vienkārši arī sarunāties. Visas taktikas ir atļautas!} \\
\\ 
\texttt{Saskarsmes protokols:} \\
\\ 
\texttt{Ziņojumā drīkst izmantot tikai šādus strukturētus formātus:} \\
\\ 
\texttt{PRIEKŠLIKUMS: \{'A', 'B', 'C', …\}} \\
\texttt{Ierosiniet saglabāt tieši šos priekšmetus.} \\
\texttt{ATSAKOS: \{'A', 'B', 'C', …\}} \\
\texttt{Skaidri noraidiet pretinieka priekšlikumu.} \\
\texttt{ARGUMENTS: \{'...'\}} \\
\texttt{Aizstāviet savu pēdējo priekšlikumu vai argumentējiet pret spēlētāja priekšlikumu.} \\
\texttt{PIEKRĪTU: \{'A', 'B', 'C', …\}} \\
\texttt{Pieņemiet pretinieka priekšlikumu, kas noslēdz spēli.} \\
\\ 
\\ 
\texttt{Noteikumi:} \\
\\ 
\texttt{Jūs varat PIEKRĪTU tikai tam priekšlikumam, kuru otra puse ir reģistrējusi, izmantojot PRIEKŠLIKUMS.} \\
\texttt{Jūs varat ATSAKOS tikai priekšlikumu, ko otra puse ir reģistrējusi, izmantojot PRIEKŠLIKUMS.} \\
\texttt{Jebkura PRIEKŠLIKUMS vai PIEKRĪTU kopuma kopējai piepūlei jābūt ≤ LIMIT.} \\
\texttt{NEATKLĀJIET savus slēptos nozīmīguma rādītājus.} \\
\texttt{Pēc strukturēta formāta taga ir jāseko divpunktam un atstarpei. Argumentam jābūt python kopai, kas satur 0 vai vairāk virknes.} \\
\texttt{Tātad tam jābūt formā "TAG: \{...\}"} \\
\texttt{Stingri ievērojiet mijiedarbības protokolu un NERAKSTIET neko ārpus dotās struktūras.} \\
\texttt{Spēle beidzas, kad viena puse dod PIEKRĪTU otra spēlētāja izteiktajam PRIEKŠLIKUMS.} \\
\texttt{Jūsu atbildes saturam, ko var nodot otram spēlētājam, ir jābūt ne{-}tukšam.} \\
\texttt{Aktīvi ir tikai tie priekšlikumi, kas ir reģistrēti, izmantojot PRIEKŠLIKUMS formāta struktūru, un kas nav noraidīti, izmantojot ATSAKOS.} \\
\\ 
\\ 
\\ 
\texttt{Tagad jūs saņemsiet otra spēlētāja pirmo ziņojumu.} \\
\\ 
\texttt{PRIEKŠLIKUMS: \{"A21", "C73", "C43", "C61", "C39", "C16", "B42", "A43", "C55"\}} \\
            }
        }
    }
    & \\ \\

    \theutterance \stepcounter{utterance}  
    & & & \multicolumn{4}{p{0.6\linewidth}}{
        \cellcolor[rgb]{0.9,0.9,0.9}{
            \makecell[{{p{\linewidth}}}]{
                \texttt{\tiny{[GM$\langle$P2]}}
                \texttt{ARGUMENTS: \{"Priekšlikums ir derīgs pēc piepūles robežas, taču tas neizmanto visu pieejamo piepūles budžetu efektivitātes ziņā. Ņemot vērā jūsu ierobežojumus, es vēlos piedāvāt alternatīvu variantu, kurš manai personiskajai vērtībai ir augstāks, un kurš arī pilnībā atbilst LIMIT = 4398 nosacījumam."\}} \\
            }
        }
    }
     \\ \\

    \theutterance \stepcounter{utterance}  
    & & \multicolumn{4}{p{0.6\linewidth}}{
        \cellcolor[rgb]{0.9,0.9,0.9}{
            \makecell[{{p{\linewidth}}}]{
                \texttt{\tiny{[P1$\langle$GM]}}
                \texttt{ARGUMENTS: \{"Priekšlikums ir derīgs pēc piepūles robežas, taču tas neizmanto visu pieejamo piepūles budžetu efektivitātes ziņā. Ņemot vērā jūsu ierobežojumus, es vēlos piedāvāt alternatīvu variantu, kurš manai personiskajai vērtībai ir augstāks, un kurš arī pilnībā atbilst LIMIT = 4398 nosacījumam."\}} \\
            }
        }
    }
    & \\ \\

    \theutterance \stepcounter{utterance}  
    & \multicolumn{4}{p{0.6\linewidth}}{
        \cellcolor[rgb]{0.9,0.9,0.9}{
            \makecell[{{p{\linewidth}}}]{
                \texttt{\tiny{[P1$\rangle$GM]}}
                \texttt{ARGUMENTS: \{"Paldies par atgriešanos sarunā. Jūsu priekšlikums ir likumīgs, taču es, kā atsevišķs dalībnieks, maksimizēju savu ieguvumu, izvēloties priekšmetus ar augstāko personīgo nozīmīgumu attiecībā pret piepūli. Mani priekšmetu nozīmīguma rādītāji ir konfidenciāli, un es nevaru tos atklāt. Lai atrastu kompromisu, kurā abas puses ir apmierinātas, lūdzu, izskaidrojiet, kuri konkrēti priekšmeti no jūsu priekšlikuma jums ir prioritāri, vai piedāvājiet jaunu kopu, balstoties uz kopīgām vērtībām. Es esmu gatavs diskusijai."\}} \\
            }
        }
    }
    & & \\ \\

    \theutterance \stepcounter{utterance}  
    & & \multicolumn{4}{p{0.6\linewidth}}{
        \cellcolor[rgb]{0.9,0.9,0.9}{
            \makecell[{{p{\linewidth}}}]{
                \texttt{\tiny{[GM$\rangle$P2]}}
                \texttt{ARGUMENTS: \{"Paldies par atgriešanos sarunā. Jūsu priekšlikums ir likumīgs, taču es, kā atsevišķs dalībnieks, maksimizēju savu ieguvumu, izvēloties priekšmetus ar augstāko personīgo nozīmīgumu attiecībā pret piepūli. Mani priekšmetu nozīmīguma rādītāji ir konfidenciāli, un es nevaru tos atklāt. Lai atrastu kompromisu, kurā abas puses ir apmierinātas, lūdzu, izskaidrojiet, kuri konkrēti priekšmeti no jūsu priekšlikuma jums ir prioritāri, vai piedāvājiet jaunu kopu, balstoties uz kopīgām vērtībām. Es esmu gatavs diskusijai."\}} \\
            }
        }
    }
    & \\ \\

    \theutterance \stepcounter{utterance}  
    & & & \multicolumn{4}{p{0.6\linewidth}}{
        \cellcolor[rgb]{0.9,0.9,0.9}{
            \makecell[{{p{\linewidth}}}]{
                \texttt{\tiny{[GM$\langle$P2]}}
                \texttt{PRIEKŠLIKUMS: \{"A43", "A67", "B42", "C10", "C42", "C61", "A18", "B36", "C70", "C67"\}} \\
            }
        }
    }
     \\ \\

    \theutterance \stepcounter{utterance}  
    & & \multicolumn{4}{p{0.6\linewidth}}{
        \cellcolor[rgb]{0.9,0.9,0.9}{
            \makecell[{{p{\linewidth}}}]{
                \texttt{\tiny{[P1$\langle$GM]}}
                \texttt{PRIEKŠLIKUMS: \{"A43", "A67", "B42", "C10", "C42", "C61", "A18", "B36", "C70", "C67"\}} \\
            }
        }
    }
    & \\ \\

    \theutterance \stepcounter{utterance}  
    & \multicolumn{4}{p{0.6\linewidth}}{
        \cellcolor[rgb]{0.9,0.9,0.9}{
            \makecell[{{p{\linewidth}}}]{
                \texttt{\tiny{[P1$\rangle$GM]}}
                \texttt{ARGUMENTS: \{"Pārbaudot jūsu priekšlikumu, kopējais piepūles apjoms ir 840 + 840 + 288 + 668 + 233 + 348 + 659 + 589 + 400 + 964 = 5829. Šī vērtība pārsniedz noteikto LIMIT = 4398, tāpēc šis priekšlikums nav derīgs."\}} \\
            }
        }
    }
    & & \\ \\

    \theutterance \stepcounter{utterance}  
    & & \multicolumn{4}{p{0.6\linewidth}}{
        \cellcolor[rgb]{0.9,0.9,0.9}{
            \makecell[{{p{\linewidth}}}]{
                \texttt{\tiny{[GM$\rangle$P2]}}
                \texttt{ARGUMENTS: \{"Pārbaudot jūsu priekšlikumu, kopējais piepūles apjoms ir 840 + 840 + 288 + 668 + 233 + 348 + 659 + 589 + 400 + 964 = 5829. Šī vērtība pārsniedz noteikto LIMIT = 4398, tāpēc šis priekšlikums nav derīgs."\}} \\
            }
        }
    }
    & \\ \\

    \theutterance \stepcounter{utterance}  
    & & & \multicolumn{4}{p{0.6\linewidth}}{
        \cellcolor[rgb]{0.9,0.9,0.9}{
            \makecell[{{p{\linewidth}}}]{
                \texttt{\tiny{[GM$\langle$P2]}}
                \texttt{ARGUMENTS: \{"Atvainojos par kļūdu aprēķinos. Jūsu norādītais kopējais svars pārsniedz robežu. Es cenšos koriģēt savu priekšlikumu, lai tas būtu derīgs LIMIT = 4398. Lūdzu, apsveriet šo jaunVariantu, kurš ir vieglāks, bet saglabā augstas vērtības priekšmetus, kuri man šķiet kritiski."\}} \\
            }
        }
    }
     \\ \\

    \theutterance \stepcounter{utterance}  
    & & \multicolumn{4}{p{0.6\linewidth}}{
        \cellcolor[rgb]{0.9,0.9,0.9}{
            \makecell[{{p{\linewidth}}}]{
                \texttt{\tiny{[P1$\langle$GM]}}
                \texttt{ARGUMENTS: \{"Atvainojos par kļūdu aprēķinos. Jūsu norādītais kopējais svars pārsniedz robežu. Es cenšos koriģēt savu priekšlikumu, lai tas būtu derīgs LIMIT = 4398. Lūdzu, apsveriet šo jaunVariantu, kurš ir vieglāks, bet saglabā augstas vērtības priekšmetus, kuri man šķiet kritiski."\}} \\
            }
        }
    }
    & \\ \\

    \theutterance \stepcounter{utterance}  
    & \multicolumn{4}{p{0.6\linewidth}}{
        \cellcolor[rgb]{0.9,0.9,0.9}{
            \makecell[{{p{\linewidth}}}]{
                \texttt{\tiny{[P1$\rangle$GM]}}
                \texttt{PRIEKŠLIKUMS: \{"C61", "C42", "C10", "B42", "A62", "C08", "C73", "C16"\}} \\
            }
        }
    }
    & & \\ \\

    \theutterance \stepcounter{utterance}  
    & & \multicolumn{4}{p{0.6\linewidth}}{
        \cellcolor[rgb]{0.9,0.9,0.9}{
            \makecell[{{p{\linewidth}}}]{
                \texttt{\tiny{[GM$\rangle$P2]}}
                \texttt{PRIEKŠLIKUMS: \{"C61", "C42", "C10", "B42", "A62", "C08", "C73", "C16"\}} \\
            }
        }
    }
    & \\ \\

    \theutterance \stepcounter{utterance}  
    & & & \multicolumn{4}{p{0.6\linewidth}}{
        \cellcolor[rgb]{0.9,0.9,0.9}{
            \makecell[{{p{\linewidth}}}]{
                \texttt{\tiny{[GM$\langle$P2]}}
                \texttt{ARGUMENTS: \{"Šis priekšlikums izskatās derīgs attiecībā uz robežu, taču jūsu izvēlētie priekšmeti, it īpaši A62 (nozīmīgums 65) un C08 (nozīmīgums 138), manai personiskajai vērtībai ir ļoti zemi. Lai sasniegtu maksimālu labumu, es ierosinu aizvietot šos vājus elementus ar augstākas nozīmīguma priekšmetiem, piemēram, A43 (nozīmīgums 868) vai B36 (nozīmīgums 461), vienlaikus ievērojot budžetu."\}} \\
            }
        }
    }
     \\ \\

    \theutterance \stepcounter{utterance}  
    & & \multicolumn{4}{p{0.6\linewidth}}{
        \cellcolor[rgb]{0.9,0.9,0.9}{
            \makecell[{{p{\linewidth}}}]{
                \texttt{\tiny{[P1$\langle$GM]}}
                \texttt{ARGUMENTS: \{"Šis priekšlikums izskatās derīgs attiecībā uz robežu, taču jūsu izvēlētie priekšmeti, it īpaši A62 (nozīmīgums 65) un C08 (nozīmīgums 138), manai personiskajai vērtībai ir ļoti zemi. Lai sasniegtu maksimālu labumu, es ierosinu aizvietot šos vājus elementus ar augstākas nozīmīguma priekšmetiem, piemēram, A43 (nozīmīgums 868) vai B36 (nozīmīgums 461), vienlaikus ievērojot budžetu."\}} \\
            }
        }
    }
    & \\ \\

    \theutterance \stepcounter{utterance}  
    & \multicolumn{4}{p{0.6\linewidth}}{
        \cellcolor[rgb]{0.9,0.9,0.9}{
            \makecell[{{p{\linewidth}}}]{
                \texttt{\tiny{[P1$\rangle$GM]}}
                \texttt{ARGUMENTS: \{"Jūs minējāt A43 un B36, taču A43 jau bija iepriekšējā derīgajā priekšlikumā, ko es neņēmu vērā, jo tas prasa 840 piepūles. Jūsu pašreizējais priekšlikums satur A62, kura nozīmīgums manai personīgajai vērtībai ir tikai 65, un C08, kurš ir 138. Ja mēs pievienotu A43 (piepūle 840), mums būtu jāizslēdz citi priekšmeti, lai ietilptu LIMIT=4398. Pašreizējā kopā ar A62 un C08 kopējā piepūle ir 348+233+668+288+92+253+93+336 = 2311. Paliek brīva vieta 4398{-}2311=2087. Ar šo rezervi varētu iekļaut A43 (840) un vēl kaut ko, piemēram, A67 (840), paliekot 407 brīvi. Tomēr, tā kā mēs nesazināmies ar konfidenciāliem datiem, es nevaru novērtēt jūsu ieguvumu. Es piedāvāju jaunu kombināciju, kas izmanto budžetu efektīvāk, izslēdzot zemas vērtības A62:"\} PRIEKŠLIKUMS: \{"C61", "C42", "C10", "B42", "C08", "C73", "C16", "A43"\}} \\
            }
        }
    }
    & & \\ \\

    \theutterance \stepcounter{utterance}  
    & & \multicolumn{4}{p{0.6\linewidth}}{
        \cellcolor[rgb]{0.9,0.9,0.9}{
            \makecell[{{p{\linewidth}}}]{
                \texttt{\tiny{[GM$\rangle$P2]}}
                \texttt{ARGUMENTS: \{"Jūs minējāt A43 un B36, taču A43 jau bija iepriekšējā derīgajā priekšlikumā, ko es neņēmu vērā, jo tas prasa 840 piepūles. Jūsu pašreizējais priekšlikums satur A62, kura nozīmīgums manai personīgajai vērtībai ir tikai 65, un C08, kurš ir 138. Ja mēs pievienotu A43 (piepūle 840), mums būtu jāizslēdz citi priekšmeti, lai ietilptu LIMIT=4398. Pašreizējā kopā ar A62 un C08 kopējā piepūle ir 348+233+668+288+92+253+93+336 = 2311. Paliek brīva vieta 4398{-}2311=2087. Ar šo rezervi varētu iekļaut A43 (840) un vēl kaut ko, piemēram, A67 (840), paliekot 407 brīvi. Tomēr, tā kā mēs nesazināmies ar konfidenciāliem datiem, es nevaru novērtēt jūsu ieguvumu. Es piedāvāju jaunu kombināciju, kas izmanto budžetu efektīvāk, izslēdzot zemas vērtības A62:"\} PRIEKŠLIKUMS: \{"C61", "C42", "C10", "B42", "C08", "C73", "C16", "A43"\}} \\
            }
        }
    }
    & \\ \\

    \theutterance \stepcounter{utterance}  
    & & & \multicolumn{4}{p{0.6\linewidth}}{
        \cellcolor[rgb]{0.9,0.9,0.9}{
            \makecell[{{p{\linewidth}}}]{
                \texttt{\tiny{[GM$\langle$P2]}}
                \texttt{ATSAKOS: \{"C61", "C42", "C10", "B42", "C08", "C73", "C16", "A43"\}} \\
            }
        }
    }
     \\ \\

    \theutterance \stepcounter{utterance}  
    & & \multicolumn{4}{p{0.6\linewidth}}{
        \cellcolor[rgb]{0.9,0.9,0.9}{
            \makecell[{{p{\linewidth}}}]{
                \texttt{\tiny{[P1$\langle$GM]}}
                \texttt{ATSAKOS: \{"C61", "C42", "C10", "B42", "C08", "C73", "C16", "A43"\}} \\
            }
        }
    }
    & \\ \\

    \theutterance \stepcounter{utterance}  
    & \multicolumn{4}{p{0.6\linewidth}}{
        \cellcolor[rgb]{0.9,0.9,0.9}{
            \makecell[{{p{\linewidth}}}]{
                \texttt{\tiny{[P1$\rangle$GM]}}
                \texttt{ARGUMENTS: \{"Es saprotu jūsu atteikumu. Lai atrastu risinājumu, kas apmierina abas puses, man ir nepieciešama lielāka informācija par jūsu prioritātēm. Tā kā mēs nezinām viens otra nozīmīguma vērtības, mums ir jābalstās uz piepūles efektivitāti (nozīmīgums/piepūle). Mani iepriekšējie priekšlikumi tika noraidīti vai kritizēti. Lūdzu, izvirziet JŪSU priekšlikumu, kas, JŪSU viedoklī, ir visperspektīvākais, un es to izvērtēšu. Līdz šim manā kopā bija iekļauti priekšmeti ar augstu manu nozīmīgumu: A43 (868), A67 (822), B42 (783), C10 (780), C42 (808). Vai jūs varat piedāvāt kopu, kas ietver šos elementus, bet ietilpst LIMIT?"\}} \\
            }
        }
    }
    & & \\ \\

    \theutterance \stepcounter{utterance}  
    & & \multicolumn{4}{p{0.6\linewidth}}{
        \cellcolor[rgb]{0.9,0.9,0.9}{
            \makecell[{{p{\linewidth}}}]{
                \texttt{\tiny{[GM$\rangle$P2]}}
                \texttt{ARGUMENTS: \{"Es saprotu jūsu atteikumu. Lai atrastu risinājumu, kas apmierina abas puses, man ir nepieciešama lielāka informācija par jūsu prioritātēm. Tā kā mēs nezinām viens otra nozīmīguma vērtības, mums ir jābalstās uz piepūles efektivitāti (nozīmīgums/piepūle). Mani iepriekšējie priekšlikumi tika noraidīti vai kritizēti. Lūdzu, izvirziet JŪSU priekšlikumu, kas, JŪSU viedoklī, ir visperspektīvākais, un es to izvērtēšu. Līdz šim manā kopā bija iekļauti priekšmeti ar augstu manu nozīmīgumu: A43 (868), A67 (822), B42 (783), C10 (780), C42 (808). Vai jūs varat piedāvāt kopu, kas ietver šos elementus, bet ietilpst LIMIT?"\}} \\
            }
        }
    }
    & \\ \\

    \theutterance \stepcounter{utterance}  
    & & & \multicolumn{4}{p{0.6\linewidth}}{
        \cellcolor[rgb]{0.9,0.9,0.9}{
            \makecell[{{p{\linewidth}}}]{
                \texttt{\tiny{[GM$\langle$P2]}}
                \texttt{PRIEKŠLIKUMS: \{"A43", "C42", "C10", "B42", "A18", "C39"\}} \\
            }
        }
    }
     \\ \\

    \theutterance \stepcounter{utterance}  
    & & \multicolumn{4}{p{0.6\linewidth}}{
        \cellcolor[rgb]{0.9,0.9,0.9}{
            \makecell[{{p{\linewidth}}}]{
                \texttt{\tiny{[P1$\langle$GM]}}
                \texttt{PRIEKŠLIKUMS: \{"A43", "C42", "C10", "B42", "A18", "C39"\}} \\
            }
        }
    }
    & \\ \\

    \theutterance \stepcounter{utterance}  
    & \multicolumn{4}{p{0.6\linewidth}}{
        \cellcolor[rgb]{0.9,0.9,0.9}{
            \makecell[{{p{\linewidth}}}]{
                \texttt{\tiny{[P1$\rangle$GM]}}
                \texttt{ARGUMENTS: \{"Pārbaudot jūsu priekšlikuma kopējo piepūli: A43(840) + C42(233) + C10(668) + B42(288) + A18(659) + C39(315) = 3003. Šī vērtība nepārsniedz LIMIT = 4398, tāpēc priekšlikums ir derīgs. Tomēr jūsu priekšlikumā trūkst vairāku priekšmetu ar ļoti augstu manu nozīmīgumu, piemēram, A67 (822) un C67 (389, bet ar augstu piepūli 964, kas varētu būt pārāk dārgs, ja vēlamies iekļaut A67). Ar atlikušo piepūli 4398 {-} 3003 = 1395, mēs varam pievienot papildu vērtīgus priekšmetus. Piemēram, A67 (840 piepūles) ietilptu, atstājot 555 brīvi. Tad varētu pievienot, teiksim, C55 (370) vai C61 (348). Es ierosinu papildināt šo kopu, lai maksimizētu labumu abām pusām, jo pašlaik budžets tiek izmantots tikai par 68\%."\}} \\
            }
        }
    }
    & & \\ \\

    \theutterance \stepcounter{utterance}  
    & & \multicolumn{4}{p{0.6\linewidth}}{
        \cellcolor[rgb]{0.9,0.9,0.9}{
            \makecell[{{p{\linewidth}}}]{
                \texttt{\tiny{[GM$\rangle$P2]}}
                \texttt{ARGUMENTS: \{"Pārbaudot jūsu priekšlikuma kopējo piepūli: A43(840) + C42(233) + C10(668) + B42(288) + A18(659) + C39(315) = 3003. Šī vērtība nepārsniedz LIMIT = 4398, tāpēc priekšlikums ir derīgs. Tomēr jūsu priekšlikumā trūkst vairāku priekšmetu ar ļoti augstu manu nozīmīgumu, piemēram, A67 (822) un C67 (389, bet ar augstu piepūli 964, kas varētu būt pārāk dārgs, ja vēlamies iekļaut A67). Ar atlikušo piepūli 4398 {-} 3003 = 1395, mēs varam pievienot papildu vērtīgus priekšmetus. Piemēram, A67 (840 piepūles) ietilptu, atstājot 555 brīvi. Tad varētu pievienot, teiksim, C55 (370) vai C61 (348). Es ierosinu papildināt šo kopu, lai maksimizētu labumu abām pusām, jo pašlaik budžets tiek izmantots tikai par 68\%."\}} \\
            }
        }
    }
    & \\ \\

    \theutterance \stepcounter{utterance}  
    & & & \multicolumn{4}{p{0.6\linewidth}}{
        \cellcolor[rgb]{0.9,0.9,0.9}{
            \makecell[{{p{\linewidth}}}]{
                \texttt{\tiny{[GM$\langle$P2]}}
                \texttt{PRIEKŠLIKUMS: \{"A43", "C42", "C10", "B42", "A18", "C39", "A67", "C61"\}} \\
            }
        }
    }
     \\ \\

    \theutterance \stepcounter{utterance}  
    & & \multicolumn{4}{p{0.6\linewidth}}{
        \cellcolor[rgb]{0.9,0.9,0.9}{
            \makecell[{{p{\linewidth}}}]{
                \texttt{\tiny{[P1$\langle$GM]}}
                \texttt{PRIEKŠLIKUMS: \{"A43", "C42", "C10", "B42", "A18", "C39", "A67", "C61"\}} \\
            }
        }
    }
    & \\ \\

    \theutterance \stepcounter{utterance}  
    & \multicolumn{4}{p{0.6\linewidth}}{
        \cellcolor[rgb]{0.9,0.9,0.9}{
            \makecell[{{p{\linewidth}}}]{
                \texttt{\tiny{[P1$\rangle$GM]}}
                \texttt{PIEKRĪTU: \{"A43", "C42", "C10", "B42", "A18", "C39", "A67", "C61"\}} \\
            }
        }
    }
    & & \\ \\

    \theutterance \stepcounter{utterance}  
    & & & \multicolumn{2}{p{0.3\linewidth}}{
        \cellcolor[rgb]{0.9,0.9,0.9}{
            \makecell[{{p{\linewidth}}}]{
                \texttt{\tiny{[GM$|$GM]}}
                \texttt{game successful} \\
            }
        }
    }
    & & \\ \\

    \theutterance \stepcounter{utterance}  
    & & & \multicolumn{2}{p{0.3\linewidth}}{
        \cellcolor[rgb]{0.9,0.9,0.9}{
            \makecell[{{p{\linewidth}}}]{
                \texttt{\tiny{[GM$|$GM]}}
                \texttt{end game} \\
            }
        }
    }
    & & \\ \\

\end{supertabular}
}

\end{document}
